\documentclass[aspectratio=169,11pt,t,xcolor=table]{beamer}
\usepackage{slidestyle}
\externaldocument[H-]{handout}
% Lesson 3 lost-mark points (from the material plus teaching observation), numbered once
\setcounter{lostmark}{0}
\lostmark{tpcopy}{Copying the shape: the \tturns of $f$ become \kw{$x$-intercepts} of $f'$, not \tturns of $f'$.}
\lostmark{sign}{Wrong side of the axis: where $f$ goes down, $f'$ is \kw{below} the $x$-axis.}
\lostmark{degree}{Wrong type of graph: a parabola's $f'$ is a straight line; a cubic's $f'$ is a parabola.}
\lostmark{steep}{The steepest point of $f$ not matched: it gives the \tturn of $f'$.}
\lostmark{roots}{Reverse sketch: the zeros of $f'$ are the \tturns of $f$, not where $f$ crosses the axis.}
\lostmark{marks}{Key $x$-values not lined up with the given graph, or the new graph not labelled.}

\newlength\figunit
\def\figxstep{2}\def\figystep{4}% coarser tick labels on screen
\begin{document}

% ------------------------------------------------------------------
\begin{frame}{Lesson 3 \textperiodcentered\ Sketching Gradient Functions}\relax
\vskip6mm
{\sEmph\bfseries Learning intention\par}\vskip2mm
Sketch the graph of $f'(x)$ from the graph of $f(x)$, and sketch a possible $f(x)$ from the graph of $f'(x)$.
\vskip6mm
{\sEmph\bfseries Success looks like\par}\vskip2mm
\begin{itemize}
\item I put the zeros of $f'$ under the \tturns of $f$.
\item I put $f'$ on the correct side of the axis in every interval.
\item I can go backwards from $f'$ to a possible $f$.
\end{itemize}
\note{Teaching line, not a question slide.\par
Materials: the practice sheet starts with a sketch page; grids 1 to 6 are for the three practice rounds.\par
Timing: Do Now 10 min, three rounds about 20 min each, practice sheet 30 min, wrap-up 5 min.}
\end{frame}

% ------------------------------------------------------------------
\begin{frame}{Do Now}\relax
\begin{columns}[T]
\begin{column}{0.52\textwidth}
\qn{1} Find the \tturn of $y=x^{2}-6x+1$ and state its nature.
\vskip5mm
\qn{2} For which $x$ is $f(x)=x^{3}-3x^{2}$ decreasing?
\vskip5mm
\qn{3} Is the gradient at A, B, C and D positive, negative or zero?
\end{column}
\begin{column}{0.46\textwidth}\centering
\setlength\figunit{12mm}
% y = x^3 - 3x with A (x=-1.5), B (x=-1), C (x=0.5), D (x=1.7); gradients 3.75, 0, -2.25, 5.67.
\begin{tikzpicture}[x=\figunit, y=0.4\figunit, line cap=round]
  \draw[-{Stealth[length=2mm]}, line width=0.6pt] (-2.4,0) -- (2.6,0) node[right, inner sep=1pt] {$x$};
  \draw[-{Stealth[length=2mm]}, line width=0.6pt] (0,-3.5) -- (0,4) node[above, inner sep=1pt] {$y$};
  \begin{scope}\clip (-2.4,-3.5) rectangle (2.6,3.8); \curve{-2.2}{2.2}{\x*\x*\x-3*\x}\end{scope}
  \foreach \p/\n/\a in {-1.5/A/north west, -1/B/south, 0.5/C/south west, 1.7/D/north west}{
    \fill (\p,{\p*\p*\p-3*\p}) circle (2.2pt);
    \node[anchor=\a, inner sep=3pt] at (\p,{\p*\p*\p-3*\p}) {\n};}
\end{tikzpicture}

\end{column}
\end{columns}
\sfoot{Lesson 2 and prerequisites}
\note{Q1 $(3,\,-8)$, minimum\par
Q2 $0<x<2$\par
Q3 A positive; B zero; C negative; D positive\par
Short solution:\par
Q1 $2x-6=0$; $y=9-18+1$; $y''=2>0$\par
Q2 $f'(x)=3x^{2}-6x=3x(x-2)<0$ between the roots\par
Q3 $f'(x)=3x^{2}-3$: $3.75$, $0$, $-2.25$, $5.67$\par
Watch for:\par
Q3 reading the sign of $y$ instead of the direction of the curve (A is below the axis but going up)\par
Diagnostic:\par
Q1 or Q2 wrong: Lesson 2 needs a recap before round 2\par
Q3 wrong: the student confuses height and gradient, the exact confusion this lesson fixes; spend longer on the one idea}
\end{frame}

% ------------------------------------------------------------------
\begin{frame}{The one idea}\relax
\fbx{\text{height of } f' \;=\; \text{steepness of } f}
\vskip2mm
\renewcommand{\arraystretch}{1.35}
\centering
\begin{tabular}{@{}p{62mm}p{62mm}@{}}
\toprule
graph of $f$ & graph of $f'$\\
\midrule
\tturn & \\
going up & \\
going down & \\
steepest point & \\
\bottomrule
\end{tabular}
\sfoot{\hp{table}}
\note{Row 1 on the $x$-axis ($f'=0$)\par
Row 2 above the $x$-axis\par
Row 3 below the $x$-axis\par
Row 4 turning point of $f'$\par
Short solution:\par
fill the table with the student before any graph; keep it on the whiteboard all lesson\par
Watch for:\par
L1: students expect the same hump in the same place; say “flat means zero”}
\end{frame}

% ------------------------------------------------------------------
\begin{frame}{Lines and parabolas}\relax
A straight line: $f'$ is a \kw{horizontal line}. A parabola: $f'$ is a \kw{straight line} through the $x$-value of its \tturn.
\vskip2mm
\qn{1} Sketch $f'(x)$ for this parabola, on the axes on the right.
\vskip2mm
\centering
\setlength\figunit{7mm}
% f(x) = x^2 - 2x - 3 = (x-1)^2 - 4: minimum at x = 1.  Caller sets \figunit.
\sketch[xmin=-2,xmax=4,ymin=-5,ymax=6,unit=\figunit,yunit=0.5\figunit,xstep=\figxstep,ystep=\figystep,name={$f(x)$},
  curve={\x*\x-2*\x-3},dmin=-2,dmax=4,
  extra={\draw[dashed, line width=0.6pt] (1,-5) -- (1,6);}]
\hspace{12mm}% f'(x) = 2x - 2: crosses the x-axis at x = 1.  Curve drawn only when \showfp is defined.
\sketch[xmin=-2,xmax=4,ymin=-6,ymax=6,unit=\figunit,yunit=0.5\figunit,xstep=\figxstep,ystep=\figystep,name={$f'(x)$},
  curve={2*\x-2},show=\ifdefined\showfp 1\else 0\fi,dmin=-2,dmax=4,
  extra={\draw[dashed, line width=0.6pt] (1,-6) -- (1,6);}]

\sfoot{\hp{parab}}
\note{Q1 straight line crossing the $x$-axis at $x=1$, below the axis for $x<1$, above for $x>1$\par
Short solution:\par
Q1 here $f(x)=x^{2}-2x-3$, so $f'(x)=2x-2$: through $(1,\,0)$ and $(0,\,-2)$\par
Watch for:\par
L3: drawing $f'$ as a curve\par
the dashed line at $x=1$ is the ruler line: the zero of $f'$ sits on it}
\end{frame}

% ------------------------------------------------------------------
\begin{frame}{Practice}\relax
Sketch $f'(x)$ on the sketch page, grids 1 and 2.
\vskip2mm
\begin{columns}[T]
\begin{column}{0.48\textwidth}\centering
\qn{1}\\[1mm]
\sketch[xmin=-3,xmax=3,ymin=-4,ymax=6,unit=8mm,yunit=4.5mm,xstep=2,ystep=2,name={$f(x)$},curve={-2*\x+3},dmin=-3,dmax=3]
\end{column}
\begin{column}{0.48\textwidth}\centering
\qn{2}\\[1mm]
\sketch[xmin=-1,xmax=5,ymin=-5,ymax=5,unit=8mm,yunit=4.3mm,xstep=2,ystep=2,name={$f(x)$},curve={-\x*\x+4*\x},dmin=-1,dmax=5]
\end{column}
\end{columns}
\sfoot{Sketch page}
\note{Q1 horizontal line at height $-2$\par
Q2 straight line through $(2,\,0)$, above the axis for $x<2$, below for $x>2$\par
Short solution:\par
Q1 the line falls $2$ for every $1$ across: gradient $-2$ everywhere\par
Q2 maximum at $x=2$; $f(x)=-x^{2}+4x$, so $f'(x)=-2x+4$, through $(0,\,4)$\par
Watch for:\par
Q1 drawing $f'$ as a sloping line\par
Q2 slope direction: up then down means $f'$ goes from $+$ to $-$}
\end{frame}

% ------------------------------------------------------------------
\begin{frame}{Cubics and quartics}\relax
{\sEmph\bfseries Write it in this order\par}\vskip2mm
\begin{enumerate}
\item \tturns of $f$ $\to$ zeros of $f'$ (use a ruler).
\item $f$ going up $\to$ $f'$ above the axis; going down $\to$ below.
\item Steepest point of $f$ $\to$ \tturn of $f'$.
\item Join smoothly. Check the type: cubic $\to$ parabola, quartic $\to$ cubic.
\end{enumerate}
\sfoot{\hp{cubic}}
\note{Teaching line, not a question slide.\par
Model step 1 with a ruler on the whiteboard: vertical lines from each turning point down to the new axis.}
\end{frame}

% ------------------------------------------------------------------
\begin{frame}{Worked example \textperiodcentered\ cubic}\relax
\qn{2} Sketch $f'(x)$ for $f(x)=x^{3}-3x$, on the axes on the right.
\vskip2mm
\centering
\setlength\figunit{8.5mm}
% f(x) = x^3 - 3x: maximum at x = -1, minimum at x = 1, steepest downhill at x = 0.
\sketch[xmin=-3,xmax=3,ymin=-6,ymax=6,unit=\figunit,yunit=0.45\figunit,xstep=\figxstep,ystep=\figystep,name={$f(x)$},
  curve={\x*\x*\x-3*\x},dmin=-2.6,dmax=2.6,
  extra={\draw[dashed, line width=0.6pt] (-1,-6) -- (-1,6) (1,-6) -- (1,6);}]
\hspace{12mm}% f'(x) = 3x^2 - 3: zeros at x = -1 and 1, minimum -3 at x = 0.  Curve only with \showfp.
\sketch[xmin=-3,xmax=3,ymin=-4,ymax=8,unit=\figunit,yunit=0.45\figunit,xstep=\figxstep,ystep=\figystep,name={$f'(x)$},
  curve={3*\x*\x-3},show=\ifdefined\showfp 1\else 0\fi,dmin=-2.2,dmax=2.2,
  extra={\draw[dashed, line width=0.6pt] (-1,-4) -- (-1,8) (1,-4) -- (1,8);}]

\sfoot{\hp{cubic}}
\note{Q2 up-parabola with zeros at $x=-1$ and $x=1$, lowest point at $x=0$\par
Short solution:\par
Q2 up, down, up gives above, below, above; $f'(x)=3x^{2}-3$, minimum $-3$\par
Watch for:\par
L4: the lowest point of $f'$ must be under the steepest downhill part of $f$ ($x=0$)}
\end{frame}

% ------------------------------------------------------------------
\begin{frame}{Practice}\relax
Sketch $f'(x)$ on the sketch page, grids 3 and 4.
\vskip2mm
\begin{columns}[T]
\begin{column}{0.48\textwidth}\centering
\qn{1}\\[1mm]
\sketch[xmin=-4,xmax=3,ymin=-16,ymax=24,unit=8mm,yunit=1.15mm,xstep=2,ystep=8,gy=4,name={$f(x)$},curve={\x*\x*\x+1.5*\x*\x-6*\x},dmin=-4,dmax=3]
\end{column}
\begin{column}{0.48\textwidth}\centering
\qn{2}\\[1mm]
\sketch[xmin=-3,xmax=3,ymin=-18,ymax=10,unit=8mm,yunit=1.6mm,xstep=2,ystep=8,gy=2,name={$f(x)$},curve={\x*\x*\x*\x-8*\x*\x},dmin=-3,dmax=3]
\end{column}
\end{columns}
\sfoot{Sketch page}
\note{Q1 up-parabola with zeros at $x=-2$ and $x=1$\par
Q2 cubic with zeros at $x=-2$, $0$ and $2$: below, above, below, above\par
Short solution:\par
Q1 maximum at $x=-2$, minimum at $x=1$; $f'(x)=3x^{2}+3x-6=3(x+2)(x-1)$\par
Q2 quartic with minimums at $x=\pm2$ and a maximum at $x=0$; $f'(x)=4x^{3}-16x$\par
Watch for:\par
Q1 lowest point of $f'$ at $x=-0.5$, halfway between the zeros\par
Q2 Merit: a quartic gives a cubic with three zeros}
\end{frame}

% ------------------------------------------------------------------
\begin{frame}{From $f'$ back to $f$}\relax
\renewcommand{\arraystretch}{1.35}
\centering
\begin{tabular}{@{}p{62mm}p{62mm}@{}}
\toprule
$f'$ graph & $f$ graph\\
\midrule
crosses from $+$ to $-$ & \\
crosses from $-$ to $+$ & \\
above the axis & \\
below the axis & \\
\bottomrule
\end{tabular}
\vskip4mm
\raggedright
$f'$ gives gradients only, so the height of $f$ is unknown: sketch \kw{a possible} $f$.
\sfoot{\hp{reverse}}
\note{Row 1 maximum\par
Row 2 minimum\par
Row 3 increasing\par
Row 4 decreasing\par
Short solution:\par
read each row from the $f'$ graph left to right\par
Watch for:\par
L5: the zeros of $f'$ are turning points of $f$, not its $x$-intercepts}
\end{frame}

% ------------------------------------------------------------------
\begin{frame}{Worked example \textperiodcentered\ backwards}\relax
\qn{3} Sketch a possible $f(x)$ for this $f'(x)$, on the axes on the right.
\vskip2mm
\centering
\setlength\figunit{8.5mm}
% given f'(x) = -(x+1)(x-3) = -x^2 + 2x + 3: zeros at -1 and 3.
\sketch[xmin=-2,xmax=4,ymin=-6,ymax=5,unit=\figunit,yunit=0.45\figunit,xstep=\figxstep,ystep=\figystep,name={$f'(x)$},
  curve={-\x*\x+2*\x+3},dmin=-2,dmax=4,
  extra={\draw[dashed, line width=0.6pt] (-1,-6) -- (-1,5) (3,-6) -- (3,5);}]
\hspace{12mm}% a possible f(x) = -x^3/3 + x^2 + 3x: minimum (-1, -5/3), maximum (3, 9).  Curve only with \showfp.
\sketch[xmin=-2,xmax=4,ymin=-4,ymax=10,unit=\figunit,yunit=0.35\figunit,xstep=\figxstep,ystep=\figystep,name={$f(x)$},
  curve={-\x*\x*\x/3+\x*\x+3*\x},show=\ifdefined\showfp 1\else 0\fi,dmin=-2,dmax=4,
  extra={\draw[dashed, line width=0.6pt] (-1,-4) -- (-1,10) (3,-4) -- (3,10);}]

\sfoot{\hp{reverse}}
\note{Q3 minimum at $x=-1$, maximum at $x=3$, decreasing, increasing, decreasing\par
Short solution:\par
Q3 one possible $f(x)=-\dfrac{1}{3}x^{3}+x^{2}+3x$: minimum $(-1,\,-1.67)$, maximum $(3,\,9)$\par
Watch for:\par
L5: drawing $f$ through $x=-1$ and $x=3$ on the axis}
\end{frame}

% ------------------------------------------------------------------
\begin{frame}{Practice}\relax
Sketch a possible $f(x)$ on the sketch page, grids 5 and 6.
\vskip2mm
\begin{columns}[T]
\begin{column}{0.48\textwidth}\centering
\qn{1}\\[1mm]
\sketch[xmin=-3,xmax=3,ymin=-2,ymax=5,unit=8mm,yunit=5.5mm,xstep=2,ystep=2,name={$f'(x)$},curve={3},dmin=-3,dmax=3]
\end{column}
\begin{column}{0.48\textwidth}\centering
\qn{2}\\[1mm]
\sketch[xmin=-1,xmax=5,ymin=-5,ymax=6,unit=8mm,yunit=3.5mm,xstep=2,ystep=2,name={$f'(x)$},curve={\x*\x-4*\x},dmin=-1,dmax=5]
\end{column}
\end{columns}
\sfoot{Sketch page}
\note{Q1 any straight line with gradient $3$\par
Q2 maximum at $x=0$, minimum at $x=4$\par
Short solution:\par
Q1 $f'$ is $3$ everywhere, so $f$ rises $3$ for every $1$ across, e.g. $f(x)=3x$\par
Q2 $f'$ goes $+$ to $-$ at $0$ and $-$ to $+$ at $4$; one possible $f(x)=\dfrac{1}{3}x^{3}-2x^{2}+4$\par
Watch for:\par
Q1 drawing a horizontal line (copying $f'$)\par
Q2 steepest downhill at $x=2$, under the lowest point of $f'$}
\end{frame}

% ------------------------------------------------------------------
\begin{frame}{Ways to lose marks}\relax
\begin{itemize}
\item[\lmnum{tpcopy}] \tturns of $f$ become \kw{zeros} of $f'$.
\item[\lmnum{sign}] $f$ going down $\to$ $f'$ below the axis.
\item[\lmnum{degree}] Parabola $\to$ line; cubic $\to$ parabola.
\item[\lmnum{roots}] Zeros of $f'$ are \tturns of $f$.
\item[\lmnum{marks}] Line up $x$-values with a ruler; label the graph.
\end{itemize}
\sfoot{\hp{lose}}
\note{Teaching line, not a question slide.\par
Full list L1 to L6 in the handout, with the versions that get no credit.}
\end{frame}

% ------------------------------------------------------------------
\begin{frame}{Summary}\relax
{\sSum
\begin{enumerate}
\item Height of $f'$ = steepness of $f$.
\item Turning points of $f$ $\leftrightarrow$ zeros of $f'$.
\item Up: $f'$ above the axis. Down: $f'$ below.
\item Backwards: $+$ to $-$ is a maximum; the height of $f$ is free.
\end{enumerate}}
\sfoot{\hp{idea}}
\note{Teaching line, not a question slide.\par
Ask the student to say each line back before looking at the screen.}
\end{frame}

% ------------------------------------------------------------------
\begin{frame}{Practice sheet \textperiodcentered\ 30 minutes}\relax
Work alone on the practice sheet, from Question 1.
\vskip3mm
\begin{itemize}
\item Questions 1 and 2: warm-up, about 6 minutes.
\item Questions 3 to 7: exam questions, about 24 minutes.
\item Use a ruler to line up every key $x$-value.
\end{itemize}
\sfoot{Practice sheet}
\note{Q1 up-parabola, zeros at $x=-1$ and $x=3$\par
Q2 parabola with a maximum at $x=2$\par
Q3 cubic: zeros under the three turning points; negative, positive, negative, positive\par
Q4 cubic shape: minimum under the left zero of $f'$, maximum under the right zero\par
Q5 negative quartic shape: maximum at $x=-3$, minimum at $x=-1$, maximum at $x=3.5$\par
Q6 (i) $f'=-1$, $g'=1$, $h'(x)=2x$\par
Q7 (ii) $p'=-1$ for $x<0$ and $p'=1$ for $x>0$; (iii) no\par
Short solution:\par
Q1 $f'(x)=x^{2}-2x-3=(x+1)(x-3)$\par
Q2 $f'$ is $+$ then $-$, crossing at $x=2$\par
Q5 $f'$ crosses $+$ to $-$ at $-3$ and $3.5$, $-$ to $+$ at $-1$\par
Q7 (iii) $p$ has a sharp corner at $x=0$; the gradient jumps from $-1$ to $1$, so it is not defined there, not zero\par
Watch for:\par
Q6 (i) straight lines give horizontal lines; $h'$ passes through the origin\par
Q7 open circles at $x=0$ on $p'$\par
Full working: practice answer version (tutor's working, not an official mark scheme)}
\end{frame}

% ------------------------------------------------------------------
\begin{frame}{Homework}\relax
{\sEmph\bfseries Homework sheet: 14 questions, about 60 minutes\par}
\vskip4mm
\begin{itemize}
\item Questions 1 to 11: Achieved and Merit sketches.
\item Questions 12 and 13: Merit.
\item Question 14: Extension (Excellence).
\end{itemize}
\vskip4mm
Optional extra: StudyTime Guide, ``Sketching Gradient Functions'' from page 19.
\sfoot{Homework}
\note{Teaching line, not a question slide.\par
Answers: homework answer version (tutor's working, not an official mark scheme).\par
Next lesson Do Now checks one $f$ to $f'$ sketch in words.}
\end{frame}
\end{document}
